\documentclass[11pt]{article}
\usepackage[utf8]{inputenc}
\usepackage[T1]{fontenc}
\usepackage{amsmath}
\usepackage{amsfonts}
\usepackage{amssymb}
\usepackage{geometry}
\usepackage{hyperref}

\geometry{letterpaper, margin=1in}

\title{\textbf{Problem A: Ship Wars}}
\author{Time Limit: 2.0s | Memory Limit: 256MB}
\date{}

\begin{document}

\maketitle

\section*{Description}

Hudson, Wiley, and Jayden are engaged in a heated debate about which anime has better relationships (``ships''): \textit{Bocchi the Rock!} or \textit{Oshi no Ko}. Unable to come to a consensus through discussion alone, they decide to settle the matter with a vote.

They have compiled a list of $N$ potential ships. For each ship, they have assigned a numerical score $S_i$ representing its quality (based on chemistry, lore, and fanart potential). Each ship belongs to exactly one fandom: either ``Bocchi'' or ``OshiNoKo''.

Your task is to help them calculate the total score for each fandom and determine the winner.

\begin{itemize}
    \item If the total score for ``Bocchi'' is strictly greater than the total score for ``OshiNoKo'', the winner is \textbf{Bocchi}.
    \item If the total score for ``OshiNoKo'' is strictly greater than the total score for ``Bocchi'', the winner is \textbf{OshiNoKo}.
    \item If the total scores are equal, the result is a \textbf{Tie}.
\end{itemize}

\section*{Input}

The first line contains a single integer $N$ ($1 \le N \le 1000$), the number of ships.

The next $N$ lines each describe a ship. Each line contains:
\begin{enumerate}
    \item A string $F_i$ representing the fandom, which is either ``Bocchi'' or ``OshiNoKo''.
    \item An integer $S_i$ ($1 \le S_i \le 100$), the score assigned to that ship.
\end{enumerate}

\section*{Output}

Output a single line containing the result:
\begin{itemize}
    \item ``Bocchi'' if the Bocchi fandom wins.
    \item ``OshiNoKo'' if the Oshi no Ko fandom wins.
    \item ``Tie'' if the total scores are equal.
\end{itemize}

\section*{Examples}

\subsection*{Example 1}
\textbf{Input:}
\begin{verbatim}
3
Bocchi 10
OshiNoKo 8
Bocchi 5
\end{verbatim}

\textbf{Output:}
\begin{verbatim}
Bocchi
\end{verbatim}

\textbf{Explanation:} \\
Total score for Bocchi: $10 + 5 = 15$. \\
Total score for OshiNoKo: $8$. \\
$15 > 8$, so Bocchi wins.

\subsection*{Example 2}
\textbf{Input:}
\begin{verbatim}
2
Bocchi 50
OshiNoKo 60
\end{verbatim}

\textbf{Output:}
\begin{verbatim}
OshiNoKo
\end{verbatim}

\textbf{Explanation:} \\
Total score for Bocchi: $50$. \\
Total score for OshiNoKo: $60$. \\
$60 > 50$, so OshiNoKo wins.

\subsection*{Example 3}
\textbf{Input:}
\begin{verbatim}
2
Bocchi 100
OshiNoKo 100
\end{verbatim}

\textbf{Output:}
\begin{verbatim}
Tie
\end{verbatim}

\end{document}
